% ---
% resumo em português
\begin{resumo}
\noindent
Inconsistências na classificação fiscal de produtos em sistemas de gestão empresarial (ERPs), como a Nomenclatura Comum do Mercosul (NCM) vazia, inexistente ou de outro produto, levam à rejeição da Nota Fiscal Eletrônica ou à tributação incorreta e geram demanda recorrente ao suporte técnico. Este trabalho desenvolve e avalia um protótipo de classificação fiscal NCM assistida por Modelos de Linguagem de Grande Escala (LLMs) executados localmente, integrado ao banco de dados Firebird do ERP Lux Auto, no qual o modelo apenas sugere e nenhuma alteração é gravada sem aprovação humana. Códigos inválidos são detectados por consulta determinística à tabela oficial. A sugestão combina a expansão de abreviações, uma busca lexical (SQLite FTS5 com BM25, sinônimos e fusão de rankings), uma base de conhecimento com produtos já classificados por duas lojas de autopeças, consultada por vizinhos mais próximos, e duas escolhas do modelo Qwen3 8B restritas por esquema JSON, de modo que o modelo só pode responder com códigos existentes. Códigos válidos são apontados como possivelmente incorretos quando ficam fora das três sugestões e o modelo julga o texto oficial incompatível com o produto. A avaliação usou conjuntos separados de desenvolvimento e de teste, este medido uma única vez. A detecção apontou todos os códigos vazios, malformados, inexistentes e extintos, sem falsos positivos em mais de 41 mil produtos reais; para códigos válidos, porém incorretos, teve acurácia de 86,6\% no desenvolvimento e de 88,8\% em erros reais de uma loja, em que apontou todos os 40 produtos cadastrados com um código ``coringa'', e de 86,0\% no teste, com 6,0\% de alarmes falsos. Na sugestão, a versão final acertou o código em primeiro lugar em 48,0\% e entre as três opções em 58,0\% dos 50 produtos de teste, em 13,7 segundos por produto, contra 20,0\% e 30,0\% da busca sem modelo de linguagem; nos produtos automotivos, foco do ERP, acertou 83,3\% em primeiro lugar. A meta de 75\% de acerto foi atingida apenas no recorte automotivo, e a maior parte dos erros restantes decorre da distância entre o nome comercial do produto e o texto oficial da tabela. Os resultados indicam que um LLM pequeno executado localmente é viável como assistente de classificação fiscal, desde que restrito a escolher entre candidatos verificados, apoiado em exemplos reais e subordinado à aprovação humana.

\vspace{\onelineskip}
\noindent
\textbf{Palavras-chave}: Modelos de Linguagem de Grande Escala. Classificação Fiscal. NCM. Geração Aumentada por Recuperação. \textit{Human-in-the-Loop}. ERP.
\end{resumo}

% resumo em inglês
\begin{resumo}[Abstract]
\begin{otherlanguage*}{english}
\noindent
Inconsistencies in the tax classification of products in Enterprise Resource Planning (ERP) systems, such as a Mercosur Common Nomenclature (NCM) code that is empty, nonexistent or belongs to another product, cause Brazilian electronic invoices to be rejected or taxed incorrectly and create recurring demand for technical support. This work develops and evaluates a prototype for NCM classification assisted by locally run Large Language Models (LLMs), integrated with the Firebird database of the Lux Auto ERP, in which the model only suggests and no change is written without human approval. Invalid codes are detected by a deterministic lookup in the official table. Suggestion combines abbreviation expansion, lexical retrieval (SQLite FTS5 with BM25, synonyms and rank fusion), a knowledge base of products already classified by two auto parts stores, queried by nearest neighbors, and two choices by the Qwen3 8B model constrained by a JSON schema, so the model can only answer with existing codes. Valid codes are flagged as possibly incorrect when they fall outside the three suggestions and the model judges the official text incompatible with the product. The evaluation used separate development and test sets, the latter measured only once. Detection flagged every empty, malformed, nonexistent and discontinued code, with no false positives among more than 41 thousand real products; for valid but incorrect codes, it reached 86.6\% accuracy in development and 88.8\% on real errors from one store, where it flagged all 40 products registered with a ``catch-all'' code, and 86.0\% on the test set, with a 6.0\% false alarm rate. For suggestion, the final version ranked the correct code first for 48.0\% and among the three options for 58.0\% of the 50 test products, in 13.7 seconds per product, against 20.0\% and 30.0\% for retrieval without a language model; for automotive products, the focus of the ERP, it ranked the correct code first for 83.3\%. The 75\% accuracy target was met only for automotive products, and most remaining errors come from the gap between commercial product names and the official table text. The results indicate that a small, locally run LLM is viable as a tax classification assistant, provided it is restricted to choosing among verified candidates, supported by real examples and subordinate to human approval.

\vspace{\onelineskip}
\noindent
\textbf{Keywords}: Large Language Models. Tax Classification. NCM. Retrieval-Augmented Generation. Human-in-the-Loop. ERP.
\end{otherlanguage*}
\end{resumo}
